% Source printed pp. 47-49.
\section{Grassmann 函子, Hilbert 函子与 Quot 函子}\label{src:2.1}
\begin{epigraph}
至于观念, 实体, 抽象和超验之物, 我始终无法让他们的脑袋理解哪怕最微小的概念.
\epigraphauthor{Jonathan Swift, ``格列佛游记''}
\end{epigraph}
\subsection{主要定理的陈述}\label{src:2.1.1-subsection}
\begin{theorem}[Grassmann 概形的射影性]\label{src:2.1.1}
\index{Grassmannian scheme@Grassmann 概形\quad Grassmannian scheme|hyperpage}
设 $V$ 为 Noether 概形 $S$ 上秩为 $n$ 的向量丛. 对整数 $0<q<n$, 函子
\[
\begin{aligned}
\Gr(q,V):\Sch&\longrightarrow\Sets,\\
(T\to S)&\longmapsto
\{\text{$T$ 上秩为 $q$ 的向量丛商 $V_T\twoheadrightarrow Q$}\}
\end{aligned}
\]
可由 $S$ 上的射影概形表出.
\end{theorem}
当 $S=\operatorname{Spec}\mathbb Z$ 且 $V=\mathcal O_{\operatorname{Spec}\mathbb Z}^{\oplus n}$ 时, 函子 $\Gr(q,V)$ 等同于例~\bookref{1.3.17}~中引入的 Grassmann 函子 $\Gr(q,n)$, 且
\[
\Gr(q,\mathcal O_S^{\oplus n})\cong\Gr(q,n)\times_{\mathbb Z}S.
\]
\Needspace{5\baselineskip}
\begin{theorem}[Hilbert 概形的射影性]\label{src:2.1.2}
\index{Hilbert scheme@Hilbert 概形\quad Hilbert scheme|hyperpage}
对每个 Noether 概形 $S$ 和每个多项式 $P\in\mathbb Q[z]$, 函子
\[
\begin{aligned}
\Hilb^P(\mathbb P_S^n/S):\Sch/S&\longrightarrow\Sets,\\
(T\to S)&\longmapsto
\left\{\begin{array}{l}
\text{在 $T$ 上平坦且有限表现的闭子概形 $Z\subseteq\mathbb P_T^n$, }\\
\text{使每个 $t\in T$ 对应的 $Z_t\subseteq\mathbb P_{\kappa(t)}^n$}\\
\text{都具有 Hilbert 多项式 $P$}
\end{array}\right\}
\end{aligned}
\]
可由 $S$ 上的射影概形表出.
\end{theorem}
注意, $\Hilb^P(\mathbb P_S^n/S)=\Hilb^P(\mathbb P_{\mathbb Z}^n/\mathbb Z)\times_{\mathbb Z}S$.
\begin{theorem}[Quot 概形的射影性]\label{src:2.1.3}
\index{Quot scheme@Quot 概形\quad Quot scheme|hyperpage}
设 $S$ 为 Noether 概形, $F$ 为 $\mathbb P_S^n$ 上的凝聚层, 且对某些 $l,r$, 它可以写成 $\mathcal O_{\mathbb P_S^n}(-l)^{\oplus r}$ 的商. 对每个多项式 $P\in\mathbb Q[z]$, 函子
\[
\begin{aligned}
\Quot^P(F/\mathbb P_S^n/S):\Sch/S&\longrightarrow\Sets,\\
(T\to S)&\longmapsto
\left\{\begin{array}{l}
\text{$\mathbb P_T^n$ 上的拟凝聚且有限表现的商 $F_T\twoheadrightarrow Q$, }\\
\text{其中 $Q$ 在 $T$ 上平坦, 且每个 $t\in T$ 对应的}\\
\text{$Q_t:=Q|_{\mathbb P_{\kappa(t)}^n}$ 都具有 Hilbert 多项式 $P$}
\end{array}\right\}
\end{aligned}
\]
可由 $S$ 上的射影概形表出.
\end{theorem}
特别地, $\Hilb^P(\mathbb P_S^n/S)$ 和 $\Quot^P(F/\mathbb P_S^n/S)$ 都是拟紧的!
\begin{caution}\label{src:2.1.4}
当基概形 $S$ 明确时, 我们常将记号简写为 $\Hilb^P(\mathbb P_S^n)$ 和 $\Quot^P(F/\mathbb P_S^n)$. 我们也将滥用记号, 用 $\Gr(q,V)$, $\Hilb^P(\mathbb P_S^n)$ 和 $\Quot^P(F/\mathbb P_S^n)$ 同时表示函子及表出它的概形.
\end{caution}
\begin{remark}\label{src:2.1.5}
\begin{enumerate}[(1)]
\item Grassmann 概形和 Hilbert 概形都是 Quot 概形的特例: 
\[
\Gr(q,V)\cong\Quot^P(V/\mathbb P_S^0),
\qquad
\Hilb^P(\mathbb P_S^n)=\Quot^P(\mathcal O_{\mathbb P_S^n}/\mathbb P_S^n),
\]
其中第一式中的 $P(z)=q$ 为常数多项式.
\item 在上述 Grassmann 函子和 Quot 函子的定义中, 若两个商 $p:F_T\twoheadrightarrow Q$ 和 $p':F_T\twoheadrightarrow Q'$ 的核作为 $F_T$ 的子层相等, 即 $\ker(p)=\ker(p')$, 则将它们等同; 等价地, 存在同构 $\alpha:Q\xrightarrow{\sim}Q'$, 使 $p=p'\circ\alpha$. 在 Hilbert 函子中, $X_T$ 的两个子概形在作为子概形相等时被等同(等价地, 其理想层作为 $\mathcal O_{X_T}$ 的子层相等).
\item 当 $T$ 为 Noether 概形时, $\Hilb^P(\mathbb P_S^n)$ 和 $\Quot^P(F/\mathbb P_S^n)$ 的定义中, $Z$ 在 $T$ 上有限表现及 $Q$ 有限表现的条件都是多余的.
\item 若不固定 Hilbert 多项式 $P$, 上述定义扩展为函子 $\Hilb(\mathbb P_S^n)$ 和 $\Quot(F/\mathbb P_S^n)$, 它们可由局部有限型概形表出. 有分解
\[
\Hilb(\mathbb P_S^n)=\coprod_P\Hilb^P(\mathbb P_S^n),
\qquad
\Quot(F/\mathbb P_S^n)=\coprod_P\Quot^P(F/\mathbb P_S^n).
\]
这是半连续性定理(\bookref{A.6.4})的推论, 因为 $\mathbb P_T^n$ 上在 $T$ 上平坦的层 $Q$ 的 Hilbert 多项式局部常值. 还要回忆, 当基概形 $T$ 约化时, 商层 $Q$ 在 $T$ 上平坦等价于 Hilbert 多项式局部常值(命题~\bookref{A.2.4}).
\item 若 $S$ 为仿射概形, 则 $\mathbb P_S^n$ 上每个凝聚层 $F$ 都可以对某些 $l,r$ 写成 $\mathcal O_{\mathbb P_S^n}(-l)^{\oplus r}$ 的商.
\item Hilbert 概形和 Quot 概形的可表性与射影性在更一般的情形下成立. 例如, 对闭子概形 $X\subseteq\mathbb P_S^n$, 可以用同样的方式陈述并证明函子 $\Hilb^P(X)$ 和 $\Quot^P(F/X)$ 的主要定理. 更一般地, 若 $X$ 为 $S$ 上某个向量丛 $E$ 的射影化 $\mathbb P(E)$ 的闭子概形, 即 $X\to S$ 强射影, 这些定理也成立. 推广将在 \S\bookref{2.4.4-subsection} 讨论, Chow 概形等变体将在 \S\bookref{2.4.5-subsection} 讨论.
\end{enumerate}
\end{remark}
\textit{本章路线.} 在 \S\bookref{2.2} 中, 我们将 $\Gr(q,V)$ 的射影性归约为特例
$\Gr(q,n)=\Gr(q,\mathcal O_{\operatorname{Spec}\mathbb Z}^{\oplus n})$, 
并通过函子化的 Plücker 嵌入
$\Gr(q,n)\to\mathbb P(\bigwedge^q\mathcal O_{\operatorname{Spec}\mathbb Z}^{\oplus n})$
证明它可由射影概形表出: 在概形 $S$ 上, 商 $\mathcal O_S^{\oplus n}\to Q$ 被映至线丛商 $\bigwedge^q\mathcal O_S^{\oplus n}\to\bigwedge^qQ$.
在 \S\bookref{2.3} 中, 我们引入 Castelnuovo--Mumford 正则性, 并建立 Mumford 的正则性有界性结果(\bookref{2.3.8}).
归约到 $F=\mathcal O_{\mathbb P^n}(-l)^{\oplus r}$ 的情形后, 上同调与基变换表明: 当 $d\gg0$ 时, 沿 $\pi:\mathbb P_S^n\to S$ 的正像 $\pi_*F(d)$ 是向量丛, 且其构造与基变换交换. 正则性的有界性进一步表明, 定义在 $S$-概形 $T$ 上的函子态射
\begin{equation}\label{src:2.1.6}
\begin{aligned}
\Quot^P(F/\mathbb P_S^n)&\longrightarrow\Gr(P(d),\pi_*F(d)),\\
[F_t\twoheadrightarrow Q]&\longmapsto
[\pi_{T,*}(F_T(d))\to\pi_{T,*}(Q(d))]
\end{aligned}
\end{equation}
是良定义的. 注意, 对 $F=\mathcal O_{\mathbb P_S^n}$ 时的 $\Hilb^P(\mathbb P_S^n)$, 有
$\pi_*(\mathcal O_{\mathbb P_S^n}(d))\cong\operatorname{Sym}^d\mathcal O_S^{\oplus n+1}$, 
因此~\eqref{src:2.1.6}~成为
\[
\Hilb^P(\mathbb P_S^n)\longrightarrow
\Gr(P(d),\operatorname{Sym}^d\mathcal O_S^{\oplus n+1}).
\]
然后我们证明~\eqref{src:2.1.6}~可由局部闭浸入表出(命题~\bookref{2.4.1}).
由于 $\Gr(q,\pi_*F(d))$ 可由 $S$ 上的射影概形表出, 这已给出 $\Quot^P(F/\mathbb P_S^n)$ 的可表性与拟射影性.
最后, 通过检验赋值判别法, 我们证明 $\Quot^P(F/\mathbb P_S^n)$ 在 $S$ 上固有(命题~\bookref{2.4.2}), 从而推出它在 $S$ 上射影.
本章遵循以下优秀论述: \cite[\S$14$--$15$]{Mum66a}, \cite[\S$2$]{AK80}, \cite[\S$1$]{Kol96}, \cite[\S$1.8$]{Laz04a}~和~\cite{Nit05}.

\textit{历史评注.} Grothendieck 在~\cite[定理 $3.2$]{FGAIV}~中建立了 Hilbert 概形和 Quot 概形的射影性.
我们的论述大体遵循 Grothendieck 的原始策略, 同时采用 Mumford 的简化来建立 Hilbert 概形和 Quot 概形的有界性(等价地, 有限型性).
有界性是证明中最困难的部分之一, 代数几何中几乎每个模空间的有界性论证最终都依赖于 Hilb 或 Quot 的有界性.
Grothendieck 的有界性方法依赖于 Chow 关于 Chow 概形的有界性结果, 后者参数化固定次数的约化纯维子概形.
在~\cite{Mum66a}~中, Mumford 引入了凝聚层的正则性---现在称为 Castelnuovo--Mumford 正则性---并证明: 对充分大的整数 $m$, 每个具有固定 Hilbert 多项式的子层 $F\subset\mathcal O_{\mathbb P^n}^{\oplus r}$ 都是 $m$-正则的(定理~\bookref{2.3.8}).
Mumford 利用这一结果构造曲面上曲线的 Hilbert 概形, 但他的论证同样适用于构造 Quot 概形.
